\documentclass[12pt]{article}
\usepackage{amsmath,amssymb}
\begin{document}
\section*{Research notes}

Problem: uniqueness of invariant probability for Dirichlet stochastic heat
equation on $[0,1]$ with additive space-time white noise and distributional
drift $\delta_0(u)$, under the ABLM interpretation/regularisation by
$G_{1/n}$.

Key interpretation:
\[
\delta_0(r)=-\frac12(-2{\bf1}_{r>0})',
\]
so the equation is a singular gradient SPDE with bounded potential
$f(r)=-2H_+(r)$.  For smooth approximation $B_n' = G_{1/n}$ with
$B_n(r)=\int_{-\infty}^rG_{1/n}(s)\,ds$, the invariant measure is
\[
\pi_n(dv)=Z_n^{-1}\exp\left(2\int_0^1 B_n(v(x))\,dx\right)\gamma(dv),
\]
where $\gamma$ is Brownian bridge law (OU invariant measure).  As
$n\to\infty$, $B_n\to H_+$ away from zero and Brownian bridge has zero
Lebesgue measure level set at 0 a.s., so
\[
\pi_n\Rightarrow
\pi(dv)=Z^{-1}\exp\left(2\int_0^1{\bf1}_{v(x)>0}\,dx\right)\gamma(dv).
\]
Density is bounded between positive constants, so $\pi\sim\gamma$.

Dirichlet form:
\[
\mathcal E(F,G)=\frac12\int\langle DF,DG\rangle_{L^2(0,1)}\,d\pi.
\]
Bounebache--Zambotti construct the skew heat equation for bounded BV
potential $f$ and prove the associated Markov process has invariant Gibbs
measure $\nu\propto e^{-\int f(u)}\gamma$, with resolvent kernels absolutely
continuous w.r.t. $\nu$.  For $f=-2H_+$, this matches the present drift.
Assumed ABLM weak uniqueness identifies this process with the ABLM mild
solution, so $P_t$ is symmetric in $\pi$ and has resolvent AC.

Uniqueness mechanism:
1. If $\eta$ is invariant and $R_\lambda(x,\cdot)\ll\pi$ for all $x$, then
$\eta=\lambda\eta R_\lambda\ll\pi$.
2. Let $h=d\eta/d\pi$. Symmetry gives $P_t h=h$ in $L^1(\pi)$.
3. For $h_a=h\wedge a$, concavity gives $h_a=(P_th)\wedge a\ge P_t h_a$;
integrals equal, hence fixed. Since $h_a\in L^2$, spectral gap forces
$h_a$ constant. For all $a$, this implies $h$ constant, hence $\eta=\pi$.
4. Spectral gap follows from Gaussian Poincare for $\gamma$ and bounded
density $d\pi/d\gamma$.

Potential subtlety: BZ absolute continuity is formulated for the
Dirichlet-form semigroup/resolvent; the answer interprets the ABLM solution
as the unique Markov version of this construction. If a critic insists on
using only the literal assumptions of pointwise weak convergence of
regularised solutions, those assumptions alone do not by themselves imply
resolvent absolute continuity or symmetry.
\end{document}
